\section{Threat Model}
\label{sec:threat-model}

We assume TrackOne is deployed in adversarial and failure-prone environments (remote sites, intermittent access, mixed trust on backhaul).
The key objective is that the canonical ledger artifacts (\texttt{day.bin} + Merkle metadata + \ac{OTS} proofs) remain \emph{verifiable} and \emph{non-equivocable} even if upstream networks, ingest inputs, or operational tooling degrade.

We categorize threats and failure modes along the following axes.
Each threat is paired with explicit TrackOne countermeasures and verification hooks.

\subsection{Requirement-Level Threats}
\label{subsec:threat-model}

\begin{description}[style=nextline]
  \item[T-1 Replay and ledger pollution] Attackers or misbehaving devices may replay, duplicate, or reorder frames to confuse operators and attempt to pollute the ledger.
  \textbf{Countered by:} FR-2, FR-3, FR-11, NFR-2; ledger vs.
  audit separation; and tests in \texttt{test\_replay\_*.py}.

  \item[T-2 Timestamp and calendar manipulation] \ac{OTS} calendars (public or stationary) may be unavailable, inconsistent, or malicious, leading to incomplete, delayed, or misleading proofs.
  \textbf{Countered by:} FR-6, FR-10, NFR-1, NFR-7, RR-1--RR-3; stationary calendar fallbacks; and the weekly ratchet job that upgrades pending proofs.

  \item[T-3 Commitment-encoding drift (transport vs. commitment bytes)] Ambiguity about which bytes are hashed (wire/transport encoding vs.
  canonical commitment encoding) can create silent integrity splits: roots that verify in one toolchain but cannot be reproduced by auditors.
  \textbf{Countered by:} \ac{ADR}-034; a deterministic TrackOne Canonical \ac{CBOR} profile for commitment bytes (distinct from Postcard transport bytes); and stability tests in \texttt{crates/trackone-core/src/cbor.rs}.

  \item[T-4 Cross-implementation divergence (Rust/Python/ISA)] Different gateway implementations (Rust vs.
  Python), or different CPU targets (Cortex-M vs.
  RISC-V), may disagree on Fact encoding, Merkle roots, or \ac{AEAD} semantics.
  \textbf{Countered by:} FR-12, FR-13, NFR-8; golden-corpus tests; and TrackOne-core determinism tests across targets.

  \item[T-5 Key compromise and operator-plane misuse] Theft of configuration, weak operational practices, or unauthorized operator actions could allow:
  (a) forging frames, (b) impersonating devices, or (c) suppressing/rewriting non-anchored artifacts.
  \textbf{Countered by:} C-5 (no secret logging), FR-8/FR-9 (fail-safe config), separation of duties (RR-4), immutable day artifacts once written, and explicit operator action logging (R-5).

  \item[T-6 Backhaul interception and traffic analysis] Gateways may use untrusted uplinks (public Wi-Fi/cell), enabling interception, tampering, and traffic analysis between gateway and master server.
  \textbf{Countered by:} authenticated transport hardening (e.g., a to-be-deployed WireGuard backhaul) as an operational control; and TrackOne's end-to-end artifact integrity (Merkle + \ac{OTS}) which remains verifiable even if transport confidentiality fails.
  WireGuard mitigates backhaul interception/tampering but does not protect against RF-layer attacks that occur before the gateway.

  \item[T-7 Time semantic confusion (phenomenon vs. ingest/result time)] If time fields are ambiguous, mixed, or emitted in inconsistent formats/timezones, downstream projections (SensorThings/Arc\ac{GIS}) can misrepresent events, and auditors can misinterpret the measurement window.
  \textbf{Countered by:} SA-2; explicit time-field mapping rules; and the requirement that all operator-facing timestamps be expressed as \ac{RFC}3339 (\ac{UTC}) while retaining the raw integer epoch seconds in Facts.

  \item[T-8 Identifier ambiguity (PodId formatting)] If device identifiers have inconsistent textual renderings (endianness, truncation, or mixed hex conventions), operators may mis-attribute telemetry and incident response may target the wrong pod.
  \textbf{Countered by:} fixed-width \texttt{PodId} (8 bytes) in the Fact schema; stable \texttt{PodId} conversion from legacy 32-bit IDs; and consistent hex rendering rules in tooling.

  \item[T-9 Crypto agility] ``Harvest-now, decrypt-later'' is a risk for any confidentiality-protected
  provisioning messages captured today.
  \textbf{Countered by:} suite negotiation and transcript binding in provisioning designs (hybrid-ready), while allowing pods to remain X25519-only in early deployments. \textbf{Downgrade resistance:} any negotiated \texttt{kex\_suite} MUST be authenticated and bound to the provisioning transcript.
  This threatens confidentiality, not ledger integrity; telemetry is dominated by radio airtime, not by suite choice.

  \item[T-10 Denial-of-service and resource exhaustion] Attackers may attempt to reduce availability or consume duty-cycle budget via RF jamming, replay floods, bogus-device floods, malformed frame floods, or deliberate triggering of rejection paths (log noise).
  \textbf{Countered by:} FR-4 (audit logging), NFR-2 (ledger vs.\ audit separation), and operational controls (rate limiting, allowlists, monitoring).
  DoS affects timeliness and availability; it does not invalidate the integrity of anchored artifacts.
\end{description}

\textbf{Key lifecycle note (rotation / revocation / quarantine):}
Track One treats key lifecycle as an operational control rather than a fully automated \ac{PKI} system.
When keys are rotated, the gateway updates the device configuration and resets or reinitializes the per-device session state (e.g., replay window state) as part of the provisioning procedure.
When a pod is suspected compromised, operators may revoke/quarantine it by removing its keys from the active configuration and refusing future frames from that device identifier; such frames may still be retained as audit noise but MUST NOT be admitted into \texttt{day.bin} or the Merkle set.
Automated fleet-wide rotation and long-running revocation workflows are deferred to later tracks, but the boundary is explicit here to avoid implying that cryptography alone provides fleet safety.

A more detailed attacker model and mapping to countermeasures is given in Section~\ref{sec:security}.
